\input{estilo}
% Propuesta: los cinco shards numerados y a la misma distancia. El frontend
% manda m₁ al 3 (al azar), que lo copia al 2 y al 4; debajo de cada copia, lo
% que guarda: el mensaje y la lista de dónde están las tres. Agregado del
% texto: los números y lo que guarda cada copia.
\begin{document}
\begin{figura}\begin{tikzpicture}
  \newcommand\msj[1]{\node[draw=segundo, fill=segundo!14, text=segundo, minimum size=4.6mm,
    inner sep=0pt, font=\scriptsize] at #1 {\var{m}\textsubscript{1}};}
  \foreach \i in {1,...,5} {
    \pic at ({1.9*(\i-1)},-0.55) {base};
    \node[etiq, text=suave] at ({1.9*(\i-1)},-0.95) {\i};
  }
  \draw[msg] (3.8,2.0) -- node[cod, right] {send(\var{m}\textsubscript{1})} (3.8,0.75);
  \foreach \x in {1.9, 3.8, 5.7} {\msj{(\x,-0.05)}}
  \draw[msg, draw=segundo] (3.5,0.05) to[bend right=25] (2.2,0.05);
  \draw[msg, draw=segundo] (4.1,0.05) to[bend left=25] (5.4,0.05);
  % lo que guarda cada copia
  \foreach \x in {1.9, 3.8, 5.7}
    \node[cod, text=segundo, align=center] at (\x,-1.55) {\var{m}\textsubscript{1},\\{}[2, 3, 4]};
\end{tikzpicture}\end{figura}
\end{document}
